\section*{الفصل \ref{ch:b3:petrological-thermodynamics} --- الديناميكا الحرارية للصخور}

\begin{solution}{exo:b3:petrological-thermodynamics:1}
$\Delta_rH = -3029.3 - 910.7 + 3935.0 = \qty{-5.0}{kJ}$; $\Delta_rS = 133.5 +
41.5 - 207.4 = \qty{-32.4}{J.K^{-1}}$; $\Delta_rV = 60.40 + 22.69 - 100.07 =
\qty{-16.98}{cm^3} = \qty{-1.698}{J.bar^{-1}}$. الميل $32.4/1.698 =
\qty{19.1}{bar/K}$. At \qty{873}{K}: $P = 1 + (873.15\times(-32.4) +
5000)/(-1.698) = \qty{13.7}{kbar}$.
\end{solution}

\begin{solution}{exo:b3:petrological-thermodynamics:2}
$c = 2$ و$p = 3 > c$: القاعدة منتهكة. لا تتعايش الأطوار الثلاثة إلا على
خطّ أحادي التغيّر للتفاعل \ce{Mg2SiO4 + SiO2 <=> 2 MgSiO3}؛ وصخرة مأخوذة
عشوائيًا ليست هناك، لذا يسجّل التجمّع تفاعلًا غير مكتمل، مثل
الفورستيريت المغلَّف بقشرة من الإنستاتيت التي منعته من التفاعل مع
الكوارتز.
\end{solution}

\begin{solution}{exo:b3:petrological-thermodynamics:3}
$a_{\text{غروسولار}} = 0.20^3 = 0.008$؛ $a_{\text{بيروب}} = 0.15^3 =
\num{3.4e-3}$؛ $a_{\text{أنورثيت}} = 0.40$.
\end{solution}

\begin{solution}{exo:b3:petrological-thermodynamics:4}
$T_c = 15\,000/(2\times8.314) = \qty{902}{K}$; $W = 2\times8.314\times1000
= \qty{16.6}{kJ.mol^{-1}}$.
\end{solution}

\begin{solution}{exo:b3:petrological-thermodynamics:5}
$\Delta_rG = -5432.2 + 2\times(-2088.8) - 730.2 + 3\times3478.1 =
\qty{94.3}{kJ}$؛ $P = 1 + 94\,300/6.61 = \qty{14.3}{kbar}$. الخط التجريبي
يعطي $22.80\times726.85 - 1093 = \qty{15.5}{kbar}$، \qty{1.2}{kbar}
أعلى.
\end{solution}

\begin{solution}{exo:b3:petrological-thermodynamics:6}
$RT = \qty{7483}{J.mol^{-1}}$. عند $x = 0.5$: $7483\ln0.5 + 20\,000/4 = -5187 +
5000 = \qty{-187}{J.mol^{-1}}$. عند $x = 0.111$: $7483\times(0.111\ln0.111 +
0.889\ln0.889) + 20\,000\times0.111\times0.889 = -2609 + 1974 =
\qty{-635}{J.mol^{-1}}$، والقيمة نفسها عند 0.889 بالتماثل. كميات متساوية من
الجناحين لهما التركيب الكلي 0.5 وطاقة جيبس
\qty{-635}{J.mol^{-1}}: الانفصال يخفّضها بمقدار \qty{448}{J.mol^{-1}}.
\end{solution}

\begin{solution}{exo:b3:petrological-thermodynamics:7}
$c + 2 = 5$ أطوار و5 خطوط. (د) $A + C = 2B$: $0 + 2/3 = 2\times1/3$. (أ)
$B + D = 2C$: $1/3 + 1 = 2\times2/3$. (ج) $2A + D = 3B$: $0 + 1 = 3\times1/3$.
(ب) $A + 2D = 3C$: $0 + 2 = 3\times2/3$.
\end{solution}

\begin{solution}{exo:b3:petrological-thermodynamics:8}
حسب قاعدة الرافعة، يشكّل الطور عند 0.889 نسبةً قدرها $(0.3 - 0.111)/(0.889 - 0.111) =
0.243$، \qty{24}{\percent}، والطور عند 0.111 \qty{76}{\percent}.
\end{solution}

\begin{solution}{exo:b3:petrological-thermodynamics:9}
$\partial P/\partial X_{\text{an}} = -3RT/(6.21X_{\text{an}}) =
-3\times8.314\times900/(6.21\times0.30) = \qty{-12000}{bar}$ لكل وحدة من
$X_{\text{an}}$: خطأ قدره 0.01 يحرّك الضغط بمقدار \qty{120}{bar}.
\end{solution}

\begin{solution}{exo:b3:petrological-thermodynamics:10}
$G'_{\text{خلط}} = RT\ln[x/(1 - x)] + W(1 - 2x)$ و$G''_{\text{خلط}} =
RT/[x(1 - x)] - 2W$، صفرًا حيث $x(1 - x) = RT/(2W)$. من أجل $x = 0.3$: $T =
2\times20\,000\times0.21/8.314 = \qty{1010}{K}$؛ يُبلغ \omterm{def:b3:petrological-thermodynamics:solvus}{منحنى الانحلال} عند
$20\,000\times0.4/(8.314\ln(7/3)) = \qty{1136}{K}$. بين الحدّين، البلورةُ
شبه مستقرّة: الانفصال سيخفّض طاقتها، لكنّ تقلّبًا صغيرًا
يرفعها، لذا يجب أن يتنوّى الطور الثاني.
\end{solution}

\begin{solution}{exo:b3:petrological-thermodynamics:11}
بما أنّ $G_{\text{خلط}}(x) = G_{\text{خلط}}(1 - x)$، يلامس المماسّ المشترك
المنحنى عند نقطتين متماثلتين $x$ و$1 - x$، وهو أفقي:
$G'_{\text{خلط}}(x) = 0$، وهي المعادلة. بحلّها، $T(x) = W(1 -
2x)/[R\ln((1 - x)/x)]$. مع $x = 1/2 - \varepsilon$، $1 - 2x = 2\varepsilon$
و$\ln((1/2 + \varepsilon)/(1/2 - \varepsilon)) \approx 4\varepsilon$، إذن
$T \to W/(2R)$ عندما $\varepsilon \to 0$؛ ومن أجل $x < 1/2$، يتزايد $T(x)$ كلّما اقترب $x$
من 1/2، لذا فهذه النهاية هي القيمة العظمى.
\end{solution}

\begin{solution}{exo:b3:petrological-thermodynamics:12}
$(2.0 - 0.08\times10)/0.92 = \qty{1.30}{wt\percent}$. المنغنيز يفضّل
الجارنيت: مع كمية أقل منه في \omterm{def:b3:petrological-thermodynamics:pseudosection}{التركيب الكلي الفعّال}، تنمو الحواف في
حقلٍ مختلفٍ وعلى خطوط تساوي تركيب مختلفة عمّا
كان سيتنبّأ به تحليل الصخرة الكلية.
\end{solution}

\begin{solution}{pb:b3:petrological-thermodynamics:1}
\textbf{1.} $\Delta_rH^\circ = -6640.0 + 2\times(-2593.8) - 910.7 + 3\times4234.0
= \qty{-36.3}{kJ}$; $\Delta_rS^\circ = 260.1 + 165.6 + 41.5 - 597.9 =
\qty{-130.7}{J.K^{-1}}$; $\Delta_rV = 125.28 + 88.30 + 22.69 - 302.37 =
\qty{-66.10}{cm^3} = \qty{-6.61}{J.bar^{-1}}$.

\textbf{2.} $-36\,300 + 298.15\times130.7 = \qty{+2.67}{kJ}$: موجبة، لذا
الأنورثيت مستقرّ عند السطح.

\textbf{3.} $P = 1 + (900\times(-130.7) + 36\,300)/(-6.61) = \qty{12.3}{kbar}$;
الميل $130.7/6.61 = \qty{19.8}{bar/K}$.

\textbf{4.} $1 + 81\,500/6.61 = \qty{12.3}{kbar}$ عند \qty{900}{K}؛ $1 +
119\,800/6.61 = \qty{18.1}{kbar}$ عند \qty{1200}{K}.

\textbf{5.} عند 626.85 و\qty{926.85}{\celsius} يعطي الخط التجريبي
13.2 و\qty{20.0}{kbar}: أعلى من الجداول بمقدار 0.9 و\qty{1.9}{kbar}.

\textbf{6.} $1914\times6.61 = \qty{12.7}{kJ}$. عدم اليقين المجمّع هو
$\sqrt{3.2^2 + (3\times2.0)^2 + (2\times2.0)^2 + 1.0^2} = \qty{8.0}{kJ}$: الفجوةُ
تساوي 1.6 مرّة منه، وهي ممكنة في القياسات الحرارية، وهذا هو السبب في أن الخط
يُؤخذ من التجارب.

\textbf{7.} $\Delta_rC_p$ صغيرٌ بين المواد الصلبة، ويُلغي تكاملاه جزئيًّا أحدهما الآخر:
الجداول تنقل الخط بمقدار 26 و\qty{112}{bar} فقط.

\textbf{8.} \ce{SiO2}, \ce{Al2O3}, \ce{FeO}, \ce{MgO}, \ce{MnO}, \ce{CaO},
\ce{Na2O}، \ce{K2O} و\ce{H2O}: $c = 9$. الأطوار: الجارنيت، البيوتيت،
المسكوفيت، البلاجيوكلاز، الكيانيت، الكوارتز والمائع، $p = 7$.

\textbf{9.} $F = 9 - 7 + 2 = 4$: الضغط ودرجة الحرارة ومتغيّران تركيبيّان
$p = 7 \le c = 9$: القاعدةُ محقّقةٌ، ويمكن أن تكون المجموعة المعدنية مستقرّةً
على مساحة من مستوى الضغط--درجة الحرارة.

\textbf{10.} يقع شكلان نقيّان متعدّدان للمركّب \ce{Al2SiO5} في توازن على الخط أحادي التغيّر
الكيانيت = السيليمانيت: علاقة واحدة بين $P$ و$T$،
ومقياسُ حرارةٍ يُعطي عندئذٍ الضغط.

\textbf{11.} يفرض المائعُ الجهدَ الكيميائي للماء: 8 مكوّنات حرّة
مقابل 6 أطوار صلبة، $F = 8 - 6 + 2 = 4$، و$p \le c$ لا تزال
محقّقةً.

\textbf{12.} $a_{\text{غروسولار}} = 0.10^3 = 0.001$، $a_{\text{أنورثيت}} =
0.30$; $\ln K = \ln(0.001/0.30^3) = -3.30$.

\textbf{13.} $K_D = (0.12/0.74)/(0.45/0.55) = 0.1622/0.8182 = 0.198$.

\textbf{14.} $T = (52\,635 + 2488)/(20.20 + 24.94\times1.618) = 55\,123/60.57
= \qty{910}{K}$، أو \qty{637}{\celsius}.

\textbf{15.} $P_0 = 22.80\times636.9 - 1093 = \qty{13429}{bar}$; $(3\times8.314
\times910.1/6.21)\ln(1/3) = \qty{-4016}{bar}$; $P = \qty{9.41}{kbar}$.

\textbf{16.} عند \qty{9.41}{kbar} و\qty{644.2}{\celsius}، ثمّ
\qty{9.55}{kbar}؛ وعند ذلك الضغط \qty{644.9}{\celsius}
و\qty{9.56}{kbar}، دون تغيّر بعدئذٍ: \qty{645}{\celsius}
و\qty{9.56}{kbar}.

\textbf{17.} \qty{641}{\celsius} و\qty{8.79}{kbar}. تضع الجداول خطّ الطرف النهائي
أدنى بنحو \qty{1}{kbar} (الجزء الأول)؛ لذا
يُبقى على الخط التجريبي.

\textbf{18.} عند $X_{\text{an}} = 0.35$: \qty{642}{\celsius}
و\qty{8.93}{kbar}، أقلّ بمقدار \qty{0.63}{kbar}.

\textbf{19.} عند \qty{918}{K}: $P = 1 + (918\times12.6 - 7700)/0.571 =
\qty{6.8}{kbar}$، أدنى من \qty{9.56}{kbar}: الكيانيت، كما شوهد.

\textbf{20.} $z = \num{9.56e8}/(2800\times9.81) = \qty{34.8}{km}$.

\textbf{21.} $\Delta P = (3\times8.314\times918/6.21)\ln(0.07/0.10) =
\qty{-1.3}{kbar}$: نما اللبّ على عمق أقلّ بنحو \qty{4.8}{km}. نما الجارنيت
أثناء الدفن، وتسجّل حافّته أعمق نقطة.

\textbf{22.} بلغ شيست فاردن ذروة ضغط قدرها
\textbf{\qty[text-series-to-math=true]{9.6}{kbar}} عند
\qty{645}{\celsius}، وبعمق دفن نحو \qty{35}{km}.
\end{solution}
